% TOP_PROBLEM: {"id": 10000058, "title": "Near-critical percolation limit shapes", "category_id": 19, "category": {"id": 19, "name": "probability", "display_name": "Probability", "description": "Problems involving probability theory, stochastic processes, and random structures.", "slug": "probability", "order_index": 19, "created_at": "2026-07-31T15:26:25.671Z"}, "status": "open", "problem_number": "AMR-099-0058", "set_id": 13, "statusQualification": "Specifies the intrinsic norm metric and diameter-one convention for the source’s Gromov–Hausdorff limit-shape formulation."}
% FIRST_POSTED: 2026-10-01
% ENTRY_KIND: statement_only
% UnsolvedMath ID 10000058; source catalogue code AMR-099-0058
% !TeX program = lualatex
% Definition and sources only; this file is not an LLM solution attempt.
\documentclass[11pt,a4paper]{article}
\usepackage[margin=25mm]{geometry}
\usepackage{fontspec}
\setmainfont{lmroman10-regular.otf}[BoldFont=lmroman10-bold.otf,
 ItalicFont=lmroman10-italic.otf,BoldItalicFont=lmroman10-bolditalic.otf]
\usepackage{amsmath,amssymb,mathtools}
\usepackage{unicode-math}
\setmathfont{latinmodern-math.otf}
\AtBeginDocument{\let\setminus\smallsetminus}
\usepackage{xurl,hyperref}
\hypersetup{colorlinks=true,urlcolor=blue}
\setcounter{secnumdepth}{0}
\setlength{\parindent}{0pt}
\setlength{\parskip}{5pt}
\setlength{\emergencystretch}{3em}
\begin{document}
\section*{Near-critical percolation limit shapes}\label{10000058}
{\small UnsolvedMath ID 10000058 · AMR-099-0058 · Probability · AMR Open Problem Lists}
\subsection{Definitions and mathematical statement}
For $0\le q<1/2$, independently delete each nearest-neighbor edge of $\mathbb Z^2$ with probability $q$. Let $\mathcal C_q$ be its almost-surely unique infinite open cluster, and use its chemical distance $D_q$, the minimum number of retained edges in a path. Conditional on $0\in\mathcal C_q$, the ball consists of cluster vertices at chemical distance at most $n$ from $0$.

The directional time constants define a deterministic norm $\mu_q$ on $\mathbb R^2$: for a lattice direction $x$, $D_q(\widehat0,\widehat{nx})/n\to\mu_q(x)$, with hats denoting a nearest cluster vertex, and the norm extends by homogeneity and continuity. Let $K_q=\{x:\mu_q(x)\le1\}$ be the corresponding deterministic limit shape.

View $K_q$ with metric $d_q(x,y)=\mu_q(x-y)/2$, so it has diameter one. Prove that
\[
(K_q,d_q)\xrightarrow[q\uparrow1/2]{\mathrm{GH}}
\bigl(\{x:\|x\|_2\le1/2\},\|\cdot-\cdot\|_2\bigr).
\]
Here Gromov--Hausdorff convergence means that isometric copies can be placed in common metric spaces with their Hausdorff distances tending to zero.

The two limits occur in order: first take arbitrarily large chemical-distance balls at a fixed $q$ to obtain its norm ball, then let $q$ approach the deletion threshold. Normalizing diameter removes the diverging overall distance scale and leaves the proposed recovery of rotational symmetry. The question is about graph distance inside the cluster, rather than a scaling limit of open/closed interfaces or Euclidean distances restricted to its vertices.
\subsection{Short English statement}
Do normalized chemical-distance limit shapes become Euclidean disks as square-lattice percolation approaches criticality?
\subsection{Sources}
\begin{itemize}
\item[S1] ULAM.AI and the UnsolvedMath Contributors, supplied problems.json, record 10000058 (AMR-099-0058), AMR Open Problem Lists. Catalogue identity and source transcription; CC BY 4.0.
\url{https://huggingface.co/datasets/ulamai/UnsolvedMath}.
\item[S2] Benjamini - Euclidean vs. Graph Metric, Conjecture 3.3, PDF page 4. Original source indexed by AMR-099-0058.
\url{https://www.wisdom.weizmann.ac.il/~itai/erd100.pdf#page=4}.
\item[S3] I. Benjamini, Euclidean vs. Graph Metric, primary numbered question and conventions.
\url{https://arquivo.pt/noFrame/replay/20201231041538id_/http://www.wisdom.weizmann.ac.il/~itai/erd100.pdf}.
\end{itemize}
\end{document}
